\section{Identificación del problema}
\label{sec:identificacion_problema}

A partir de los hallazgos cualitativos y cuantitativos obtenidos en el diagnóstico institucional, se establecen con precisión los siguientes tres problemas específicos que caracterizan el proceso de registro, matrícula y organización escolar en la institución:

\begin{itemize}
    \item \textbf{Ausencia de recuperación de credenciales y pérdida de cupos vocacionales:} La plataforma informática no disponía de un autoservicio para la recuperación de contraseñas ni búsqueda de usuarios asignados, lo que forzó la asistencia física masiva de estudiantes al módulo para solicitar reseteos manuales; durante esta demora presencial, las plazas de las carreras más demandadas se agotaron, impidiendo que los postulantes accedieran a la especialidad de su preferencia y obligándolos a conformarse con opciones no deseadas.
    \item \textbf{Inexistencia de un sistema para la matriculación documental posterior:} Carencia de un módulo administrativo en la gestión lectiva correspondiente para verificar, registrar y auditar en base de datos la presentación física de los requisitos obligatorios del postulante (libretas de calificaciones de 1ro a 3ro de secundaria, cédula de identidad de sus tutores, fotocopia de carnet de identidad y certificado de nacimiento), dejando trunco el ciclo formal de inscripción.
    \item \textbf{Falta de un mecanismo de asignación a cursos y paralelos por unidad educativa:} La institución carecía de una herramienta automatizada para agrupar y matricular a los estudiantes en sus cursos y paralelos específicos en función de la unidad educativa de procedencia a la que pertenecen, propiciando desequilibrios en la cantidad de alumnos por taller, errores de asignación y severos cruces de horarios entre el módulo y los colegios del distrito.
\end{itemize}
